\section{Small language models, distillation, and quantisation}\label{sec:bg-slm}
\evidence{ANALYSIS.md \S4--\S5}

I adopt the SLM survey \cite{slmsurvey2024} and on-device review \cite{ondevice2024} framing: this system is \emph{edge-only} (no cloud fallback), reported against the review's indicator template (time-to-first-token, VRAM, latency percentiles). Neither survey contains a traffic or control-loop case study, so the application gap exists by omission. The nearest on-device comparator, Octopus \cite{octopus2024}, is a 2B function-calling model at 1.1--1.7s on Android, bracketing the 0.94--1.64s median decision latency of my deployed students. The surveys corroborate a reliability floor: Qwen2.5's own report \cite{qwen25report} shows instruction-following (IFEval) collapsing from 84.1 at 72B to 27.9 at 0.5B, consistent with the 45--69\% decision accuracy of my stock 0.6B model. The Qwen3 report \cite{qwen3report} complicates this. Its 0.6B model posts respectable IFEval scores yet was the weakest decider of the Qwen3 family I served, parsing 93--100\% strict JSON while choosing wrong phases, which sharpens the point: generic instruction-following benchmarks do not measure numeric decision quality, and the technical reports never test it.

On distillation there is a tension: Orca \cite{mukherjee2023orca} argues small models need rich explanation traces, whereas my distillation is label-only. Three defences. First, the task is a bounded, closed decision schema closer to classification than open generation, for which the knowledge-distillation survey \cite{xu2024kdsurvey} endorses labelling-plus-SFT. Second, Orca itself concedes small models excel in constrained settings. Third, my student nears the teacher ceiling (97.7\% against 97.3--99.3\% inter-model label agreement on 600 cross-checked items). I concede what Orca implies: out-of-distribution generalisation and natively generated audit rationales are untested risks. The jump from 55\% to 97.7\% has precedent in distilled students beating data-scale expectations \cite{hsieh2023distilling,babyllama2023}. LoRA \cite{hu2021lora} and QLoRA \cite{dettmers2023qlora} justify rank 16, though their evidence stops above 125M parameters, and I flag the extrapolation.

Quantisation is the second and larger tension. QLoRA trains adapters on a 4-bit base and serves unmerged, whereas I train QLoRA-style, merge to bf16, then apply post-hoc int4 round-to-nearest. That is two quantisation steps, and QLoRA's authority does not transfer to the second. The literature predicts trouble there: GPTQ \cite{frantar2022gptq} shows small models are the hard case for round-to-nearest, AWQ \cite{lin2023awq} finds the gap narrowing at 4-bit with grouping, and post-hoc quantisation after fine-tuning has been measured losing 9.4\% relative even with GPTQ \cite{quantlora2024}. My contrary finding is therefore scoped narrowly: equal aggregate held-out decision accuracy, not output identity, under weight-only int4, on a near-converged narrow-schema checkpoint, measured with a discrete accuracy metric coarser than perplexity. Finally, forgetting-scaling results \cite{forgetting2024} predict general-capability drift at my very low training loss, so I carry this as a limitation (\S\ref{sec:dis-threats}), which Traffic-R1's fine-tuning-paradigm result independently motivates. What no published work supplies is per-decision reliability data for a generative model at 0.6B, int4-quantised, inside a closed control loop.
